\subsection{重数与正割序列}

命题 7. --- 假设 A 是局部环。设 s 是一个整数 ≥ 1，并且对于 \(1 \leqslant i \leqslant s\)，令 \(\delta_{i}\) 是 \textgreater{} 0 的整数，\(x_{i}\) 是 \(m_{A}^{\delta_{i}}\) 中的元素，\(\xi_{i}\) 是它在 \(m_{A}^{\delta_{i}}/m_{A}^{\delta_{i}+1}\) 中的类。假设 \((x_{1}, \ldots, x_{s})\) 是 A 的一个正割序列。令 x 表示由 \((x_{1}, \ldots, x_{s})\) 生成的 A 的理想。则有 \(e(A/x) \geqslant \delta_{1} \ldots \delta_{s} \cdot e(A)\)，当 \((\xi_{1}, \ldots, \xi_{s})\) 是 gr(A) 中的完全正割序列时取等号。

置 B = A/x，考虑形式级数

\[
\mathrm{H} _ {\mathrm{A}} = \sum_ {n \geqslant 0} \operatorname{long} (\mathfrak {m} _ {\mathrm{A}} ^ {n} / \mathfrak {m} _ {\mathrm{A}} ^ {n + 1}). \mathrm{T} ^ {n}, \quad \mathrm{H} _ {\mathrm{B}} = \sum_ {n \geqslant 0} \operatorname{long} (\mathfrak {m} _ {\mathrm{B}} ^ {n} / \mathfrak {m} _ {\mathrm{B}} ^ {n + 1}). \mathrm{T} ^ {n}
\]

并且 \(\mathrm{H}_{\mathrm{B}}^{(s)} = (1 - \mathrm{T})^{-s}\mathrm{H}_{\mathrm{B}}\)。根据第 § 4 节第 5 号的命题 6，在 \(\mathbf{Z}[[T]]\) 中有不等式

\[
\mathrm{H} _ {\mathrm{B}} ^ {(s)} \geqslant \left(\prod_ {i = 1} ^ {s} \frac {1 - \mathrm{T} ^ {\delta_ {i}}}{1 - \mathrm{T}}\right) \mathrm{H} _ {\mathrm{A}},
\]

当序列 \((\xi_{1}, \ldots, \xi_{s})\) 是完全正割序列时取等号。但是

\[
\mathrm{R(T)} = \prod_ {i = 1} ^ {s} \frac {1 - \mathrm{T} ^ {\delta_ {i}}}{1 - \mathrm{T}}
\]

这是 \(\mathbf{Z}[\mathrm{T}]\) 中的一个多项式，且 \(\mathbf{R}(1) = \delta_1 \ldots \delta_s\)。置 \(\dim(\mathbf{A}) = d\)；有 \(\dim(\mathbf{B}) = d - s\)。根据第 §4 节第 3 号的定理 2，存在 \(\mathbf{R}_{\mathbf{A}}\) 和 \(\mathbf{R}_{\mathbf{B}}\)，它们属于 \(\mathbf{Z}[\mathrm{T}, \mathrm{T}^{-1}]\)，使得

\[
\mathrm{H} _ {\mathrm{A}} = (1 - \mathrm{T}) ^ {- d} \mathrm{R} _ {\mathrm{A}} (\mathrm{T}), \quad \mathrm{R} _ {\mathrm{A}} (1) = e (\mathrm{A}),
\]

\[
\mathrm{H} _ {\mathrm{B}} = (1 - \mathrm{T}) ^ {- d + s} \mathrm{R} _ {\mathrm{B}} (\mathrm{T}), \quad \mathrm{R} _ {\mathrm{B}} (1) = e (\mathrm{A} / x).
\]

因此有

\[
(1 - \mathrm{T}) ^ {- d} \mathrm{R} _ {\mathrm{B}} (\mathrm{T}) = \mathrm{H} _ {\mathrm{B}} ^ {(s)} \geqslant \left(\prod_ {i = 1} ^ {s} \frac {1 - \mathrm{T} ^ {\delta_ {i}}}{1 - \mathrm{T}}\right) \mathrm{H} _ {\mathrm{A}} = (1 - \mathrm{T}) ^ {- d} \mathrm{R} (\mathrm{T}) \mathrm{R} _ {\mathrm{A}} (\mathrm{T}),
\]

并且当序列 \((\xi_{1},\ldots,\xi_{s})\) 是完全正割序列时取等号。根据第 § 4 节第 1 号的引理 2，结论得出。

注记。--- 反之，可以证明（参见 p. 103，练习 4），若 A 是正则的且 \(e(\mathrm{A}/\mathfrak{x})=\delta_{1}\ldots\delta_{s}\)，则序列 \((\xi_{1},\ldots,\xi_{s})\) 是完全正割序列。

例。--- 设 A 是形式幂级数环 \(k[[X_1, ..., X_n]]\)，定义在域 \(k\) 上；设 \(F_1, ..., F_s\) 是 A 的元素，令 \(\mathfrak{a}\) 为它们生成的理想，并置 \(B = A/\mathfrak{a}\)。将 \(P_1, ..., P_s \in k[X_1, ..., X_n]\) 记为级数 \(F_1, ..., F_s\) 的初始形式，并令 \(\delta_1, ..., \delta_s\) 为它们各自的次数。若序列 \(F_1, ..., F_s\) 在 A 中是正割的，则有 \(e(B) \geqslant \delta_1 ... \delta_s\)；若序列 \(P_1, ..., P_s\) 在环 \(k[X_1, ..., X_n]\) 中是完全正割的，则有 \(e(B) = \delta_1 ... \delta_s\)。

例如，考虑环 \(\mathbf{B} = k[[\mathrm{X},\mathrm{Y}]] / \mathfrak{a}\)，其中 \(\mathfrak{a}\) 由 \(\mathrm{X}^2 + \mathrm{Y}^3\) 和 \(\mathrm{X}^2 + \mathrm{Y}^4\) 生成；由上述不等式得 \(e(\mathbf{B})\geqslant 4\)；注意到 \(\mathfrak{a}\) 也由元素 \(\mathrm{X}^2 + \mathrm{Y}^3\) 和 \(\mathrm{Y}^4 -\mathrm{Y}^3\) 生成，而这些元素的初始形式序列是完全正割的，故得到 \(e(\mathbf{B}) = 6\)。
% semantic-fragments: legacy-input-seam
\relax{}
